% =========================================================
% capitulos/06_Discusion.tex
% Placeholder con las herramientas que se van a usar en la
% tesis: tablas con el estilo azul, pseudocódigo y ecuaciones
% numeradas / sin numerar de distintos tipos.
% =========================================================

\chapter{Discusión}

\label{chap:Discusión}

Los resultados obtenidos permiten analizar el comportamiento de GNLPDA desde
una perspectiva tanto geométrica como estadística, evidenciando los alcances y limitaciones
de la propuesta frente a distintos tipos de problemas de clasificación. En conjunto, los
experimentos muestran que el desempeño del método depende en gran medida de la
naturaleza del conjunto de datos, confirmando que no existe un clasificador universalmente
superior para todos los escenarios.

\section{Expansión polinomial y separabilidad de los datos}

Las proyecciones obtenidas mediante GNLPDA evidencian el efecto de incorporar
una expansión polinomial explícita antes del proceso de selección de características. En
diversos conjuntos de datos se observó que las proyecciones generadas por GNLPDA
conservan una estructura geométrica similar a la obtenida mediante LDA, aunque presentan
modificaciones suficientes para incrementar la separación entre las clases. Este
comportamiento sugiere que la expansión polinomial no sustituye la información contenida
en el espacio original, sino que incorpora términos de interacción capaces de revelar regiones
donde la separación lineal resulta más factible.

En contraste, las proyecciones obtenidas mediante KFDA presentan una geometría
considerablemente distinta. Al construir una matriz kernel y realizar la discriminación sobre
un espacio inducido implícitamente, las observaciones tienden a formar agrupamientos más
compactos dentro de cada clase, favoreciendo la separación entre grupos. Sin embargo, dicha
representación ocurre en un espacio cuya transformación no puede interpretarse
directamente, lo que dificulta identificar como están afectando las variables o que está
ocurriendo en el espacio original de las caracteristicas.

Desde esta perspectiva, GNLPDA constituye una alternativa intermedia entre los métodos
lineales y los métodos kernel. Mientras que conserva la capacidad de generar proyecciones
explícitas e interpretables, incorpora relaciones no lineales mediante la expansión polinomial
y la selección evolutiva de términos, permitiendo aproximar parte de la capacidad
discriminante de los enfoques kernel sin renunciar completamente a la trazabilidad del
modelo.

\section{Comportamiento frente al problema Small Sample Size}

El estudio realizado sobre escenarios Small Sample Size confirmó experimentalmente las limitaciones teóricas asociadas a la formulación clásica del Análisis Discriminante Lineal.
Conforme el número de variables superó ampliamente al número de observaciones, la matriz
de dispersión intra-clase perdió invertibilidad, imposibilitando la obtención de las
direcciones discriminantes mediante la formulación tradicional.

En contraste, tanto la implementación de LDA basada en descomposición en valores
singulares, disponible en el paquete MASS de R, como KFDA lograron mantener su
capacidad para generar proyecciones discriminantes bajo estas condiciones. Estos resultados
ponen de manifiesto la importancia de emplear formulaciones numéricamente robustas
cuando se trabaja con conjuntos de datos de alta dimensionalidad y tamaño de muestra
reducido, ya que permiten evitar el colapso observado en la formulación clásica sin modificar
el objetivo discriminante del método.

Asimismo, este estudio permitió ilustrar de forma visual cómo distintas
formulaciones de un mismo modelo pueden presentar comportamientos considerablemente
diferentes frente a un mismo problema, evidenciando que la implementación computacional
también desempeña un papel importante en la aplicabilidad práctica de los métodos
discriminantes.

\section{Comparación con los métodos evaluados}



La comparación experimental mostró que GNLPDA presenta un comportamiento
competitivo respecto a los métodos considerados, aunque su desempeño depende de las
características particulares del conjunto de datos.

En problemas con fronteras de decisión relativamente simples, como Palmer
Penguins o Wine, los mejores resultados fueron obtenidos por LDA y QDA, cuyos valores
de exactitud mostraron una menor variabilidad y un desempeño superior al resto de los
métodos. Estos resultados evidencian que, cuando los supuestos de los modelos clásicos se
cumplen de forma aproximada, la incorporación de mecanismos más sofisticados no
necesariamente produce mejoras significativas.

Por otra parte, el conjunto Ionosphere representó un escenario distinto debido a su elevada
dimensionalidad y al reducido número de observaciones disponibles. En este caso, QDA no
logró ejecutarse debido a las dificultades asociadas con la estimación de las matrices de
covarianza requeridas por el modelo, mientras que las formulaciones más robustas de LDA
y los métodos kernel conservaron su capacidad operativa.

En los conjuntos de datos caracterizados por fronteras de decisión claramente no lineales, los
métodos basados en kernels, particularmente la SVM con kernel gaussiano, alcanzaron los
mejores resultados de clasificación, acompañados además de una menor dispersión entre las
distintas ejecuciones. Aunque GNLPDA no logró superar este comportamiento, sí obtuvo
desempeños competitivos y estadísticamente comparables con las configuraciones basadas
en kernels polinomiales, especialmente con SVM de grado dos. Este resultado sugiere que
una expansión explícita, acompañada de un mecanismo adecuado de selección de variables,
puede aproximar la capacidad discriminante obtenida mediante kernels implícitos.

En contraste, el conjunto Pima Diabetes representó un escenario particularmente complejo
debido al elevado traslape existente entre clases. En este caso, todos los métodos evaluados
alcanzaron desempeños similares, independientemente de su complejidad. Este
comportamiento pone en evidencia que la capacidad predictiva de un clasificador no depende
únicamente del algoritmo empleado, sino también de la calidad, estructura y separabilidad
inherente de los datos disponibles.

Finalmente, los conjuntos correspondientes a los indicadores de rezago social mostraron un
comportamiento diferente al observado en los problemas anteriores. GNLPDA obtuvo un
desempeño intermedio entre los métodos clásicos y las máquinas de soporte vectorial,
mientras que KFDA presentó una disminución importante en su capacidad de clasificación.
Las proyecciones obtenidas sugieren que estos conjuntos presentan una estructura
caracterizada por clases progresivas con fronteras poco definidas, situación en la que la
representación inducida por KFDA no logró capturar adecuadamente las diferencias sutiles
entre grupos. En cambio, tanto LDA como GNLPDA conservaron una representación más
consistente de la organización gradual presente en los datos.

\section{Interpretación general de la propuesta}

En conjunto, los resultados obtenidos permiten interpretar a GNLPDA como una estrategia
que busca equilibrar capacidad predictiva, interpretabilidad y representación geométrica de
los datos. La propuesta no pretende reemplazar a los métodos clásicos cuando éstos resultan
suficientes, ni competir directamente con los enfoques kernel en aquellos escenarios donde
las transformaciones implícitas ofrecen ventajas evidentes. En cambio, plantea una
alternativa que incorpora relaciones no lineales de manera explícita, conservando la
posibilidad de identificar las variables y términos polinomiales responsables de la separación
obtenida.

Esta característica adquiere especial relevancia cuando, además del desempeño predictivo,
resulta importante comprender la estructura del problema y explicar el comportamiento del
modelo. En este sentido, GNLPDA ofrece una representación visual compacta de la
información contenida en espacios de alta dimensionalidad, permitiendo analizar la
geometría de las clases y facilitando la interpretación de las transformaciones realizadas
durante el proceso de aprendizaje.



